\documentclass[10pt,letterpaper]{article}
\usepackage[margin=0.5in]{geometry}
\usepackage{multicol}
\usepackage{amsmath,amssymb}
\usepackage{enumitem}
\usepackage{blindtext}
\usepackage{graphicx}

\setlength{\columnsep}{0.3in}
\pagestyle{empty}

\begin{document}

\begin{multicols}{2}

\textbf{Axioms of Probability}

\begin{enumerate}
    \item Non-negativity: For any event \( A \), \( P(A) \geq 0 \).
    \item Normalization: \( P(S) = 1 \), where \( S \) is the sample space.
    \item Additivity: For any two mutually exclusive events \( A \) and \( B \), \( P(A \cup B) = P(A) + P(B) \).
    This can be extended to $P(\bigcup_{i=1}^n A_i) = \sum_{i=1}^n P(A_i)$ for mutually exclusive events \( A_1, A_2, \ldots, A_n \).
\end{enumerate}

\textbf{Basic Properties of Probability:}\\
$\forall A \in S, P(A) < 1$\\
$P(A^c) = 1 - P(A)$\\
P($\emptyset$) = 0\\
$P(A \cup B) = P(A) + P(B) - P(A \cap B)$\\
$A \subseteq B \implies P(A) \leq P(B)$\\
\\
\textbf{Counting:}\\
\textbf{Permutation} --- An arrangement of objects in a specific order. $P(n, r) = \frac{n!}{(n-r)!}$\\
\textbf{Combination} --- A selection of objects without regard to the order. $C(n, r) = \frac{n!}{r!(n-r)!}$\\
\textbf{Binomial Coefficient} --- The number of ways to choose \( k \) successes in \( n \) independent Bernoulli trials. 
$\binom{n}{k} = C(n, k) = \frac{n!}{k!(n-k)!}$\\
\textbf{Number of Subsets} --- The number of subsets of a set with \( n \) elements is \( 2^n \).\\
\\
\textbf{Conditional Probability:}\\
$P(A|B) = \frac{P(A \cap B)}{P(B)}$\\
\textbf{Multiplication Rule:} --- $P(A \cap B) = P(A|B)P(B) = P(B|A)P(A)$\\
$P(\bigcap_{i=1}^n A_i) = P(A_1)P(A_2|A_1)P(A_3|A_1 \cap A_2)...P(A_n|\bigcap_{i=1}^{n-1} A_i)$\\
\textbf{Law of Total Probability:} --- If \( B_1, B_2, \ldots, B_n \) are mutually exclusive and exhaustive events, then for any event \( A \):
$P(A) = \sum_{i=1}^n P(A|B_i)P(B_i)$\\
\textbf{Bayes' Theorem:} --- \\
$P(B_j|A) = \frac{P(A|B_j)P(B_j)}{P(A)} = \frac{P(A|B_j)P(B_j)}{\sum_{i=1}^n P(A|B_i)P(B_i)}$\\
\\
\textbf{Independence:}\\
Two events \( A \) and \( B \) are independent if and only if $P(A|B) = P(A)$ and \( P(A \cap B) = P(A)P(B) \).\\
If $A$ and $B$ are independent, then $B$ and $A$, $A$ and $B^c$, $A^c$ and $B$, $A^c$ and $B^c$ are also independent.\\
\paragraph*{Random Variables}
An RV is a function that assigns a real number to each outcome in the sample space of a random experiment.\\
\textbf{Discrete RV} --- A random variable that can take on a countable number of distinct values.\\
\textbf{Continuous RV} --- A random variable that can take on any value within a given range or interval.\\
$X: S \rightarrow \mathbb{R}$ Where S is the event space and \( X \) is the random variable.\\
\textbf{PMF} --- $f(x) = P(X = x)$ (discrete)\\
$P(a \leq X \leq b) = \int_{a}^{b} f(x)dx$ (continuous)\\

\bigbreak

\textbf{CDF} --- $F(x) = P(X \leq x) = \sum_{y \leq x} P(X = y)$ (discrete)\\
$F(x) = P(X \leq x) = \int_{-\infty}^{x} f(t)dt$ (continuous)\\
$P(a \leq X \leq b) = F(b) - F(a)$\\
\textbf{Bernoulli Random Variable} --- A discrete RV that takes the value 1 with probability \( p \) and the value 0 with probability \( 1-p \).\\
\textbf{Expected Value} --- $E(X) = \mu_X = \sum_x xP(X=x)$ or $= \int_{-\infty}^{\infty} xf(x)dx$\\
$E(h(X)) = \sum_x h(x)P(X=x)$ or $E(h(X)) = \int_{-\infty}^{\infty} h(x)f(x)dx$\\
$E(X + Y) = E(X) + E(Y)$\\
$E(X^2) = \sum_x x^2p(x)$ or $= \int x^2f(x)dx$\\
\textbf{Variance} --- $\sigma_X^2 = E[(X - \mu_X)^2] = E(X^2) - (E(X))^2$ or $\sigma_X^2 = \int (x - \mu_X)^2 f(x) dx$\\
\textbf{Standard Deviation} --- $\sigma_X = \sqrt{Var(X)}$\\
\\
% \textbf{Common Discrete Distributions:}\\
% \\
% \textbf{Binomial Distribution} --- Models the number of successes in \( n \) independent Bernoulli trials, each with success probability \( p \).\\
% $X \sim Bin(n, p)$\\
% \textbf{PMF:} --- $P(X = k) = b(x; n,p) = \binom{n}{k} p^k (1-p)^{n-k}$ for $ k = 0, 1, \ldots, n $ else $0$\\
% \textbf{CDF:} --- $B(x; n,p) = P(X \leq x) = \sum_{y=0}^{\lfloor x \rfloor} b(y; n,p)$\\
% \textbf{Mean:} --- $E(X) = np$\\
% \textbf{Variance:} --- $Var(X) = np(1-p)$\\
% \\
% \textbf{Negative Binomial Distribution} --- Models the number of trials needed to achieve \( r \) successes in a sequence of independent Bernoulli trials, 
% each with success probability \( p \).\\
% $X \sim NB(r, p)$\\
% \textbf{PMF:} --- $P(X = k) = \binom{k-1}{r-1} p^r (1-p)^{k-r}$ for $ k = r, r+1, \ldots $ else $0$\\
% \textbf{CDF:} --- $NB(x; r,p) = P(X \leq x) = \sum_{y=r}^{\lfloor x \rfloor} \binom{y-1}{r-1} p^r (1-p)^{y-r}$\\
% \textbf{Mean:} --- $E(X) = \frac{r}{p}$\\
% \textbf{Variance:} --- $Var(X) = \frac{r(1-p)}{p^2}$\\
% \\
% \textbf{Geometric Distribution} --- A special case of the negative binomial distribution where \( r = 1 \). 
% It models the number of trials needed to get the first success.\\
% $X \sim Geo(p)$\\
% \textbf{PMF:} --- $P(X = k) = (1-p)^{k-1} p$ for $ k = 1, 2, \ldots $ else $0$\\
% \textbf{CDF:} --- $G(x; p) = P(X \leq x) = 1 - (1-p)^{\lfloor x \rfloor}$\\
% \textbf{Mean:} --- $E(X) = \frac{1}{p}$\\
% \textbf{Variance:} --- $Var(X) = \frac{1-p}{p^2}$\\
% \\
% \textbf{Hypergeometric Distribution} --- Models the number of successes in \( n \) draws from a finite population of size \( N \) containing \( K \) successes,
% without replacement.\\
% $X \sim Hyp(N, K, n)$\\
% \textbf{PMF:} --- $h(x; n, N, K) = P(X = k) = \frac{\binom{K}{k} \binom{N-K}{n-k}}{\binom{N}{n}}$ for $ \max(0, n-(N-K)) \leq k \leq \min(n, K) $ else $0$\\
% \textbf{CDF:} --- $H(x; n, N, K) = P(X \leq x) = \sum_{y=\max(0, n-(N-K))}^{\lfloor x \rfloor} \frac{\binom{K}{y} \binom{N-K}{n-y}}{\binom{N}{n}}$\\
% \textbf{Mean:} --- $E(X) = n \frac{K}{N}$\\
% \textbf{Variance:} --- $Var(X) = n \frac{K}{N} \frac{N-K}{N} \frac{N-n}{N-1}$\\
% \\
% \textbf{Poisson Distribution} --- Models the number of events occurring in a fixed interval of time or space, given the average rate \( \lambda \) of occurrence.\\
% Used to model the number of times a rare event occurs in a large population.\\
% $X \sim Pois(\lambda)$\\
% \textbf{PMF:} --- $P(X = k) = \frac{\lambda^k e^{-\lambda}}{k!}$ for $ k = 0, 1, 2, \ldots $ else $0$\\
% \textbf{CDF:} --- $P(X \leq x) = \sum_{y=0}^{\lfloor x \rfloor} \frac{\lambda^y e^{-\lambda}}{y!}$\\
% \textbf{Mean:} --- $E(X) = \lambda$\\
% \textbf{Variance:} --- $Var(X) = \lambda$\\
% \\

\vspace{9999pt}

\end{multicols}

\end{document}